%! AP Statistics — Unit 4, Lesson 4.1 — Lesson Plan
\documentclass[10pt]{article}
\usepackage{apstats-article}
\usepackage{apstats-boxes}
\usepackage{graphicx}
\graphicspath{{images/}}

\newcommand{\UnitNumberName}{Unit 4: Probability, Random Variables, and Probability Distributions \quad}
\newcommand{\LessonNumberName}{Lesson 4.1: Introducing Statistics --- Random and Non-Random Patterns?}

\begin{document}
\begin{center}
    {\Large\bfseries \CourseName: \SchoolYear \\
    {\small\bfseries \UnitNumberName \LessonNumberName}}
\end{center}

\begin{tcolorbox}[colback=sky, colframe=navy, boxrule=0.9pt, arc=2mm,
    left=2mm, right=2mm, top=1.5mm, bottom=1.5mm]
    \textbf{Primary Objective:} Students will recognize that streaks, clusters, and
    patterns in data can arise from pure chance, and that patterns alone do not confirm
    a non-random process (Big Idea: VAR-1).
\end{tcolorbox}

\renewcommand{\arraystretch}{1.6}
\begin{skillbox}[Priority Ideas \& Skills]{goldbox}
\begin{minipage}[t]{0.48\linewidth}
    \textbf{AP Skill 1 --- Selecting Statistical Methods}
    \begin{itemize}
        \item Identify questions suggested by patterns in data (Skill 1.A).
        \item Recognize that patterns in data do not necessarily mean the variation is
              non-random (VAR-1.F, VAR-1.F.1).
        \item Bridge from Unit 3 (data collection) to Unit 4 (modeling randomness).
    \end{itemize}
\end{minipage}
\hfill
\begin{minipage}[t]{0.48\linewidth}
    \textbf{Key Understandings}
    \begin{itemize}
        \item Random processes can produce streaks, clusters, and patterns that look
              non-random.
        \item Asking ``Could chance alone explain this?'' is a foundational statistical
              question.
        \item Patterns in data inspire questions; confirming whether those patterns are
              real requires more than observation.
        \item This lesson sets up the entire unit: probability (4.3--4.6) and random
              variables (4.7--4.12).
    \end{itemize}
\end{minipage}
\end{skillbox}

\begin{skillbox}[{Vocabulary, Concepts, \& Theorems}]{greenbox}
\begin{minipage}[t]{0.48\linewidth}
    \begin{itemize}
        \item \textbf{Random process} --- a process whose outcome cannot be predicted
              with certainty in advance
        \item \textbf{Pattern} --- an apparent regularity in observed data
        \item \textbf{Streak} --- a run of identical outcomes in sequence
    \end{itemize}
\end{minipage}
\hfill
\begin{minipage}[t]{0.48\linewidth}
    \begin{itemize}
        \item \textbf{Random variation} --- natural fluctuation produced by chance alone
        \item \textbf{Statistical question} --- a question that anticipates variability in
              the data needed to answer it
        \item \textbf{Cluster} --- an unusual concentration of outcomes in a small region
    \end{itemize}
\end{minipage}
\end{skillbox}

\begin{skillbox}[Activate Prior Knowledge \& Spiral Review]{sky}
\begin{minipage}[t]{0.48\linewidth}
    \textbf{Prior Knowledge Check (3--4 min)}
    \begin{itemize}
        \item Quick review: What is a \emph{statistical question}? (Unit 1)
        \item Connect: Unit 3 asked how data are collected; Unit 4 asks how we model
              the randomness that underlies data.
        \item Prompt: \textit{``A basketball player makes 5 free throws in a row. Is she
              on a `hot streak,' or could this happen by chance?''}
    \end{itemize}
\end{minipage}
\hfill
\begin{minipage}[t]{0.48\linewidth}
    \textbf{Connections Forward}
    \begin{itemize}
        \item Simulation to estimate probability $\to$ Lesson 4.2
        \item Formal probability rules $\to$ Lesson 4.3
        \item Conditional probability $\to$ Lessons 4.4--4.5
        \item Random variables and distributions $\to$ Lessons 4.7--4.12
        \item Quantifying how unusual a pattern is $\to$ Unit 5 (inference)
    \end{itemize}
\end{minipage}
\end{skillbox}

\begin{skillbox}[Hook \& Lesson]{sky}
\begin{minipage}[t]{0.48\linewidth}
    \textbf{Hook (5--7 min)}\\
    Display three sequences of 20 coin-flip results:
    \begin{itemize}
        \item Sequence A: HHTHTTHHHTTHHTTHTHTH
        \item Sequence B: HHHHHTTTTTHHHHHTTTTTT (10+10)
        \item Sequence C: HTHTTHTHHHTTHTHTHHTH
    \end{itemize}
    Ask: ``Which sequence looks most like real coin flips? Which looks suspicious?''\\
    Reveal: A and C were simulated; B was designed to \emph{look} suspicious. True
    random sequences often contain long runs that surprise us.
    Transition: \textit{``Our intuition about what looks random is often wrong.''}
\end{minipage}
\hfill
\begin{minipage}[t]{0.48\linewidth}
    \textbf{Lesson Flow (15--18 min)}
    \begin{enumerate}
        \item Show a sequence of basketball free-throw outcomes: a player makes 7
              consecutive shots. Ask: Does this prove she has a ``hot hand''?
        \item Key idea (VAR-1.F.1): Patterns in data do not necessarily mean the
              variation is non-random.
        \item Walk through: a run of heads in 20 coin flips is not unlikely; calculate
              intuitively (about a 12\% chance of a run of $\geq$5 in 20 flips).
        \item Distinguish the question we \emph{see} from the question we \emph{should ask}:
              ``Could this pattern have arisen by chance alone?''
        \item Preview: The rest of Unit 4 builds the probability machinery to actually
              answer that question.
    \end{enumerate}
\end{minipage}
\end{skillbox}

\begin{skillbox}[Explicit Instruction \& Active Monitoring]{sky}
\begin{minipage}[t]{0.48\linewidth}
    \textbf{Direct Instruction (10 min)}
    \begin{itemize}
        \item Show three real-world patterns: a cancer cluster in a small town, a
              winning streak in a sports team, a run of birth-year coincidences.
        \item For each: What question does the pattern suggest? Is there a non-random
              explanation? Could chance alone produce it?
        \item Emphasize VAR-1.F: pattern $\neq$ non-random; we need to model chance
              before drawing conclusions.
    \end{itemize}
\end{minipage}
\hfill
\begin{minipage}[t]{0.48\linewidth}
    \textbf{Active Monitoring}
    \begin{itemize}
        \item Listen for: ``of course that's not random!'' — redirect to ``How would you
              test that?''
        \item Cold-call: ``Give me one question this pattern suggests that we could
              actually investigate with data.''
        \item Circulate during guided notes; check that students distinguish the pattern
              they observe from the claim they are tempted to make.
    \end{itemize}
\end{minipage}
\end{skillbox}

\begin{skillbox}[Group Work \& Differentiation]{redbox}
\begin{minipage}[t]{0.48\linewidth}
    \textbf{Partner Activity (8--10 min)}\\
    Three-tier activity: students receive four scenarios (sports streaks, weather runs,
    health clusters, social media trends) and for each:
    \begin{enumerate}
        \item State the pattern you observe.
        \item Write a statistical question the pattern suggests.
        \item Decide: Could chance alone produce this pattern?
        \item What would convince you it is \emph{not} random?
    \end{enumerate}
\end{minipage}
\hfill
\begin{minipage}[t]{0.48\linewidth}
    \textbf{Differentiation}
    \begin{itemize}
        \item \textit{Support (Tier R)}: Anchor prompt --- ``A pattern is surprising.
              A non-random pattern has a \emph{cause}.'' Use guided sentence starters.
        \item \textit{On-grade (Tier A)}: Explain why human-written ``random'' sequences
              typically have too few long runs.
        \item \textit{Extension (Tier E)}: Research the ``hot hand'' fallacy (Gilovich
              et al., 1985). Summarize the statistical argument in 3 sentences.
        \item \textit{ELL}: Vocabulary card with ``random,'' ``pattern,'' and ``streak.''
    \end{itemize}
\end{minipage}
\end{skillbox}

\begin{skillbox}[Individual Work \& Assessment]{redbox}
\begin{minipage}[t]{0.48\linewidth}
    \textbf{Exit Ticket (5 min)}\\
    A student rolls a die 20 times and records: 3, 3, 3, 3, 5, 2, 1, 6, 3, 4,
    5, 2, 3, 1, 6, 4, 4, 4, 4, 4.
    \begin{enumerate}
        \item Identify two patterns in this sequence.
        \item Write one statistical question these patterns suggest.
        \item Explain why these patterns alone do not prove the die is unfair.
    \end{enumerate}
\end{minipage}
\hfill
\begin{minipage}[t]{0.48\linewidth}
    \textbf{AP-Style Check}\\
    \textit{``A researcher notices that a small town had 6 leukemia cases in one year,
    double the national average rate. Which best explains why this pattern alone does
    not confirm an environmental cause?''}
    \begin{enumerate}[label=(\Alph*)]
        \item The sample size is too large to show a real pattern.
        \item Random variation in small populations can produce clusters.
        \item Leukemia rates never vary from the national average.
        \item Six cases is not enough to calculate a probability.
    \end{enumerate}
    Answer: (B)
\end{minipage}
\end{skillbox}

\begin{skillbox}[Reinforcement \& Extension]{goldbox}
\begin{minipage}[t]{0.48\linewidth}
    \textbf{Homework / Practice}
    \begin{itemize}
        \item For four data patterns, write the statistical question each suggests and
              evaluate whether chance alone is a plausible explanation.
        \item Reflection: \textit{``Why is it dangerous to assume every pattern has a
              non-random cause?''}
    \end{itemize}
\end{minipage}
\hfill
\begin{minipage}[t]{0.48\linewidth}
    \textbf{Extension / Preview}
    \begin{itemize}
        \item \textit{Extension}: Find one news story describing a ``surprising'' streak
              or cluster. Evaluate whether the reporting correctly distinguishes chance
              from cause.
        \item \textit{Preview}: Lesson 4.2 introduces simulation as a way to model what
              chance alone would produce --- start thinking about how you would use coin
              flips to estimate probability.
    \end{itemize}
\end{minipage}
\end{skillbox}

\end{document}
